% !PRINTED 707
\weberpage{707}
%<<CONT
jede $f^{\text{te}}$ Einheitswurzel; denn wenn man im Falle 1) für $\Theta$ eine
primitive $\mu^{\text{te}}$, im Falle 2) eine primitive $\frac{1}{2}\mu^{\text{te}}$
Einheitswurzel, und im letzten Falle zugleich $\varepsilon = -1$ setzt, so ist
$\chi(n)$ eine primitive $f^{\text{te}}$ Einheitswurzel, und aus deren Potenzen
kann man alle $f^{\text{ten}}$ Einheitswurzeln herleiten.

Bedeutet daher jetzt $x$ eine Variable, so ist, wenn $n$ zum Exponenten $f$
gehört, das über sämmtliche Charaktere $\chi$ erstreckte Product
\weq{9}{\overset{\chi}{\Pi} [1 - \chi(n)\,x] = (1-x^f)^e,}
und wenn man daher $x = p^{-s}$ setzt, so ergiebt sich aus (4):
\weq{10}{\Phi(s) = \frac{1}{1-q^{-s}} \overset{\chi}{\Pi} \overset{p}{\Pi} \frac{1}{1-\chi(p)\,p^{-s}},}
wenn sich das erste Productzeichen auf alle Charaktere $\chi$, das zweite auf
alle Primzahlen $p$ erstreckt. Es ist also $\Phi(s)$ ein Product aus einer
endlichen Zahl, $\mu$, von Factoren, deren jeder ein unendliches Product ist.

Jetzt wollen wir jeden Factor eines dieser unendlichen Producte nach
steigenden Potenzen von $p^{-s}$ entwickeln. Dadurch ergiebt sich
\[
\frac{1}{1-\chi(p)\,p^{-s}} = 1 + \chi(p)\,p^{-s} + \chi(p^2)\,p^{-2s} + \chi(p^3)\,p^{-3s} + \cdots
\]
Das Product aus allen Factoren, die man hieraus erhält, wenn $\chi$
festgehalten wird, während $p$ alle von $q$ verschiedenen Primzahlen
durchläuft, ist ein Aggregat von Gliedern der Form
\[
\chi(p^k)\,\chi(p'^{k'})\,\chi(p''^{k''}) \cdots p^{-sk}\,p'^{-sk'}\,p''^{-sk''} \cdots = \chi(n)\,n^{-s},
\]
und jedes Glied dieser Form kommt in dem Producte ein- und nur einmal vor
[vgl. \S.\,192, (8)]. Danach ergiebt sich die Umformung von $\Phi(s)$ in ein
endliches Product von unendlichen Reihen:
\weq{11}{\Phi(s) = \frac{1}{1-q^{-s}} \overset{\chi}{\Pi} \overset{n}{\Sigma} \frac{\chi(n)}{n^s},}
worin sich die Summe auf alle durch $q$ nicht theilbaren Zahlen $n$ erstreckt
und das Product auf alle $\mu$ Charaktere $\chi$.

Dieser Ausdruck ist geeignet, um den Grenzwerth von $(s-1)\,\Phi(s)$ für
$s=1$ zu bestimmen.

Betrachten wir zunächst die dem Hauptcharakter $\chi = 1$ entsprechende
Summe
\[
\overset{n}{\Sigma} \frac{1}{n^s}.
\]
